\documentclass[11pt]{article}

\usepackage{amsmath,amssymb}
\usepackage{xurl}
\usepackage[hidelinks]{hyperref}
\setlength{\emergencystretch}{2em}

% Shared commands for notes in docs/paper-gaps/.
%
% Two halves.  The link commands below are project identity: OWNER/REPO is the
% placeholder that scripts/bootstrap_project.py replaces with this project's
% GitHub slug and Pages site.  Everything after them is NOTATION, and notation
% belongs to the source paper: the definitions here are an example set, and a
% project replaces them with the macros of its own paper's style file.  The one
% rule that never changes: a command defined here must mean exactly what the
% source paper's macro means, or a note silently says something else.

\newcommand{\Lean}{\textsc{Lean}}
\newcommand{\leanid}[1]{\textsf{#1}}
\newcommand{\ghissue}[1]{\href{https://github.com/OWNER/REPO/issues/#1}{GitHub issue \##1}}
\newcommand{\ghpr}[1]{\href{https://github.com/OWNER/REPO/pull/#1}{PR \##1}}
% Local issue tracker (issues/NNNN-*.md in this repository); no remote URL.
\newcommand{\localissue}[1]{issue \##1 (\path{issues/})}
\newcommand{\doclink}[1]{\href{https://OWNER.github.io/REPO/docs/#1.html}{\path{#1}}}

% --- Notation: replace with the source paper's own macros --------------------
% \F, \N, \C, \R, \tr, \ot, \eps, \E, \bx, \bu, \bv, \by

\newcommand{\F}{\mathbb{F}}
\newcommand{\N}{\mathbb{N}}
\newcommand{\C}{\mathbb{C}}
\newcommand{\R}{\mathbb{R}}
\newcommand{\tr}{\mathrm{tr}}
\newcommand{\eps}{\epsilon}
\newcommand{\ot}{\otimes}
\newcommand{\E}{\mathop{\bf E\/}}

% bold random variables
\newcommand{\bu}{\boldsymbol{u}}
\newcommand{\bv}{\boldsymbol{v}}
\newcommand{\bx}{{\boldsymbol{x}}}
\newcommand{\by}{\boldsymbol{y}}

% T1 so literal < and > in \texttt placeholders typeset as themselves
\usepackage[T1]{fontenc}

% bra-ket
\usepackage{braket}

% --- Example structured spaces ---
\newcommand{\polyfunc}[3]{\mathcal{P}(#1, #2, #3)}
\newcommand{\polymeas}[3]{\mathrm{PolyMeas}(#1, #2, #3)}
\newcommand{\polysub}[3]{\mathrm{PolySub}(#1, #2, #3)}

% Approximation relations (paper uses \approx_\eps and \simeq_\zeta)
\newcommand{\approxeps}{\approx_{\varepsilon}}
\newcommand{\simgeq}{\simeq_{\zeta}}

\renewcommand{\ghissue}[1]{%
  \href{https://github.com/Dengnifer/MIPStarRE-A/issues/#1}{GitHub issue \##1}}

\title{The Decoding Identity in the Pauli Extraction Argument}
\author{}
\date{2026-09-15}

\begin{document}
\maketitle

\section{At a glance}

\begin{itemize}
  \item \textbf{Difficulty.} Medium.  The source applies a decoder--encoder
  identity to every bounded individual-degree polynomial, although the identity
  is guaranteed for encoded polynomials and can fail outside the encoding image.
  \item \textbf{Estimated weight.} The restricted algebraic identity and the
  probabilistic non-encoding-mass estimate required algebraic and
  project-specific consistency arguments; both are now proved. No work
  remains on this support estimate. The separate source-facing consistency
  construction is also complete.
  \item \textbf{Mathlib/project split.} Evaluation of the multilinear
  interpolant is elementary polynomial algebra.  The non-encoding-mass bound
  uses project-specific consistency estimates and Schwartz--Zippel.
  \item \textbf{Key mathematical inputs.} Multilinear interpolation on the
  Boolean cube, injectivity of the low-degree encoding, measurement
  consistency, and Schwartz--Zippel.  The correction is tracked in
  \ghissue{19}.
\end{itemize}

\section{Key theorem forms}

For a word $h\colon\{0,1\}^m\to\mathbb F_q$, let $g_h$ be its multilinear
extension.  Define the chapter's decoder by
\[
  \operatorname{Dec}(g)(y)=g(y),
  \qquad y\in\{0,1\}^m.
\]
Then
\[
  \operatorname{Dec}(g_h)=h
  \quad\text{and}\quad
  \operatorname{Dec}(g_h)\mathbin{\cdot}\operatorname{ind}_m(u)=g_h(u).
\]
More generally, the second identity holds for a polynomial outcome $g$ under
the explicit hypothesis that $g$ belongs to the image of the encoding; no
converse is asserted.

For an arbitrary bounded individual-degree polynomial $g$, the Boolean values
determine its multilinear interpolant, not necessarily $g$ itself.  The
unrestricted identity
\[
  \operatorname{Dec}(g)\mathbin{\cdot}\operatorname{ind}_m(u)=g(u)
\]
is therefore false.

\section{Counterexample and source comparison}

The low-degree encoding is defined at lines~873--897 of
\path{references/qpbt-paper/04_preliminaries.tex}.  The general decoder at
lines~917--924 retains evaluations lying in a chosen subset $H$; the extraction
chapter instead specifies at lines~1419--1420 of
\path{references/qpbt-paper/14_analysis_of_the_pauli_basis_test.tex} that the
decoder returns all evaluations on the Boolean cube.  This is the specialization
$H=\mathbb F_q$, rather than the Boolean-valued shorthand used earlier
\cite[lines~873--924 and 1419--1420]{Ji2020MIPStar}.

Take $m=1$, $d\geq2$, and a field of characteristic two with more than two
elements.  For $g(x)=x^2$, the Boolean evaluations are $g(0)=0$ and $g(1)=1$.
Consequently the decoder is the word whose multilinear extension is $x$.
For any $u\notin\{0,1\}$ with $u^2\neq u$,
\[
  \operatorname{Dec}(g)\mathbin{\cdot}\operatorname{ind}_1(u)=u
  \neq u^2=g(u).
\]
Thus the equality asserted for every $g$ at lines~1483--1492 of the chapter-14
source mirror is false \cite[lines~1483--1492]{Ji2020MIPStar}.  If degrees at
least $q$ are allowed, equality as a function would not by itself prove encoding
membership either, since $x^q$ and $x$ induce the same function on $\mathbb F_q$.

\section{Required repair}

The extraction proof must split the marginal measurement into encoding and
non-encoding outcomes.  On encoding outcomes, the restricted identity above
justifies the substitution.  On the complement, the proof must bound
\[
  \sum_{g\ \mathrm{not\ an\ encoding}}
  \langle\widehat\psi,\widehat S_g^W\widehat\psi\rangle.
\]
For a supplied global polynomial-pair witness with point-consistency error
$\delta_G$, the formal proof derives an
$O(\delta_G+\sqrt\eps)+md/q$ bound. Measure the original Pauli answer $h$ and
the ancillary ideal Pauli answer $h'$ and return $g_{h+h'}$. This reference
measurement is supported on encodings. Perfect ancillary consistency and the
Pauli-basis check bound its evaluated inconsistency with the marginal.
A non-encoding polynomial differs from every encoded polynomial, so
Schwartz--Zippel bounds the probability that their evaluations coincide.
Completeness and positivity bound the non-encoding mass by the full polynomial
inconsistency. This uniform estimate holds for both player sides; it is not
inserted as an extraction hypothesis.

The printed reindexing at lines~1805--1822 replaces
$g_{\operatorname{Dec}(g_W)}$ by $g_W$, which is the invalid reverse identity
for an arbitrary polynomial outcome.  In the corrected calculation, the
constraint is
$g_h(u)=\operatorname{Dec}(g_W)\mathbin{\cdot}\operatorname{ind}_m(u)-a$.
After setting $h'=h+\operatorname{Dec}(g_W)$, the defining evaluation formula
$g_k(u)=k\mathbin{\cdot}\operatorname{ind}_m(u)$ and characteristic two turn
that constraint into $g_{h'}(u)=a$.  Thus this step uses neither decoder
linearity nor the identity $\operatorname{Dec}(g_h)=h$
\cite[lines~1805--1822]{Ji2020MIPStar}.

\section{Blueprint and formalization status}

Definition~\path{def:qld-full-field-decoder} and
Lemmas~\path{lem:qld-decoder-linearity} and
\path{lem:qld-decoder-evaluation} in
\path{blueprint/src/chapter/ch16_qpbt_extraction.tex} state the full-field
decoder, its linearity and inverse, and the evaluation identity restricted to
encoding outcomes.  Lemma~\path{lem:qld-nonencoding-mass-bound}
separately records the estimate on the remaining outcomes.

The declarations \path{decodeFq_add}, \path{decodeFq_smul}, and
\path{decodeFq_lowDegreeEncoding} in
\path{MIPStarRE/QPBT/Algebra/Decoding.lean} match the three algebraic assertions.
In the same file, \path{IsEncoding} is the equality
$g_{\operatorname{Dec}(g)}=g$ of polynomial representatives, and
\path{decodeFq_dotProduct_indicatorVec} requires an explicit hypothesis
of that predicate.  It therefore matches the corrected evaluation lemma and
does not assert the false unrestricted identity.  Finally,
\path{nonencodingMarginalMass_le} in
\path{MIPStarRE/QPBT/Extraction/Consistency.lean} states one uniform theorem for
the complementary mass bound on either player side.  It derives the bound from
the projective strategy and global polynomial-pair witness rather than adding it
as a witness field or theorem hypothesis.

The decoder calculations are proved.  \path{lowDegreeEncoding_mem_poly},
\path{decodeFq_add}, \path{decodeFq_smul}, \path{decodeFq_lowDegreeEncoding},
and \path{decodeFq_dotProduct_indicatorVec} carry complete proofs: the degree
side-condition on $g_h$, the two linearity laws, the inversion
$\operatorname{Dec}_{\mathbb F_q}(g_h)=h$, and the evaluation identity under the
explicit \path{IsEncoding} hypothesis.  The declarations
\path{tildeM_consistent_pointMeas_ofGlobalPairWitness} and
\path{tildeM_consistent_pointMeas'_ofGlobalPairWitness} prove the two specialized
non-encoding estimates used by the pulled-apart point comparisons.  Each combines
an encoding-supported reference comparison with the corresponding
Schwartz--Zippel estimate, and \path{exists_pulled_apart_consistency} applies both
results to one constructed global witness.  The uniform auxiliary theorem
\path{nonencodingMarginalMass_le} is proved at its unchanged
supplied-witness statement. Its proof uses the reference comparison and
Schwartz--Zippel from
\path{MIPStarRE/QPBT/Extraction/NonencodingSupport.lean}; it is not called by that
source-facing composition. The global witness constructor
\path{exists_globalPairWitness} and the source-facing
\path{exists_pulled_apart_consistency} are separately proved, including at
zero error. This does not remove the witness premise from the auxiliary
theorem. Issue~\#19 owns the Stage~4.2 statements and this discrepancy
record; \ghissue{47} tracks the non-encoding support theorem, completed for a
supplied witness by issue~\#517.
No unrestricted decoder identity, residual assumption, or additional witness
field was introduced.

\section{Conclusion}

The decoding identity is false for a general polynomial outcome.  Restricting
it to the encoding image is a necessary local correction.  The two specialized
non-encoding-mass estimates required for the pulled-apart point-consistency
conclusion are proved and used in the source-facing construction. The uniform
auxiliary estimate is also proved for a supplied global witness, but the
source-facing construction does not depend on it.

\bibliographystyle{alpha}
\bibliography{references}

\end{document}
